% First draft. Compile with pdflatex (Overleaf works). Swap the preamble for the venue's
% style file (rlc.sty / neurips_2026.sty / iclr2027_conference.sty) when a target is chosen.
\documentclass[11pt]{article}
\usepackage[margin=1in]{geometry}
\usepackage[utf8]{inputenc}
\usepackage[T1]{fontenc}
\usepackage{amsmath,amssymb}
\usepackage{booktabs}
\usepackage{graphicx}
\usepackage{microtype}
\usepackage{xcolor}
\usepackage[hidelinks]{hyperref}
\usepackage{caption}
\captionsetup{font=small,labelfont=bf}
\newcommand{\todo}[1]{\textcolor{red}{[TODO: #1]}}
\newcommand{\Q}{Q}
\newcommand{\dist}{d}

\title{Memorise, Interpolate, or Reason?\\
\large Diagnosing what a value network learns on an exhaustively solvable problem}
\author{Bingzhang Hu\thanks{Affiliation and e-mail. Draft v0.1.}}
\date{}

\begin{document}
\maketitle

\begin{abstract}
When a neural value function solves a combinatorial puzzle, it is rarely clear whether it has
memorised the training states, learned the structure of the problem, or something in between.
We study this question on a problem small enough to be solved exactly: the $2\times2\times2$
pocket cube, whose $3{,}674{,}160$ states and their exact distances to the solved state can be
enumerated by breadth-first search. A tabular Q-learner covers the whole space in two minutes and
serves as ground truth. We then train the same Q-learning procedure with a multilayer perceptron
and introduce four diagnostics that each rule out one explanation of its behaviour: a
seen/unseen split of the evaluation states, a \emph{symmetry-image} test that asks the network to
solve renamed copies of states it trained on, the gap between greedy and beam-search success with
the same value function, and a supervised \emph{capacity ceiling} that measures how many states
the architecture can memorise when handed the answers. On the cube the network neither memorises
individual states (seen and unseen states at the same distance are solved equally often) nor
learns the cube's symmetries (symmetric images of training states are solved no better than
random unseen states), yet its value function is informative far beyond the training
distribution: beam search of width~32 solves $90\%$ of uniformly random states where greedy
solves about $16\%$. The curriculum stalls at depth~7 because the value error ($0.95$~moves)
matches the gap between adjacent actions, not because the network lacks capacity: the same
architecture, trained supervised on all states, reaches $79\%$ greedy success in two minutes,
while it can memorise only $0.3$--$1$~million states exactly. We argue that exhaustively solvable
problems are a useful testbed for separating memorisation, interpolation and optimisation in deep
reinforcement learning, and release the code and diagnostics.
\end{abstract}

\section{Introduction}

Deep reinforcement learning (RL) has solved combinatorial puzzles such as the Rubik's cube
with a learned value function and a search on top of it \citep{mcaleer2018,agostinelli2019}.
Such results are usually reported as success rates on random instances. They leave open a
question that matters for how the method should be understood and extended: what does the value
network actually represent? Three answers are commonly given, often implicitly.
The network may have \emph{memorised} the states it saw during training and their values, so
that success on unseen states comes from those states lying a step or two from memorised
ones. It may have learned the \emph{structure} of the problem, for example the symmetries of the
puzzle, so that a solution learned for one state transfers to every equivalent state.
Or it may have learned a \emph{smooth interpolant} in the input representation that ranks
neighbouring states approximately, without either per-state memory or structural knowledge.

These explanations are hard to separate on the problems where deep RL is usually applied,
because the ground truth is unavailable: nobody knows the optimal value of a random
$3\times3\times3$ cube state, and ``unseen'' states cannot be identified without the training
trace. We therefore study a problem small enough to be solved exactly but large enough that a
network must compress it. The $2\times2\times2$ pocket cube with one cubie fixed has
$7!\cdot3^6 = 3{,}674{,}160$ reachable states, nine moves, and a maximum distance to the solved
state of~11. Breadth-first search (BFS) gives every state's exact distance in seconds, and a
bitmap of the states visited during training gives the exact seen/unseen split.

Our contributions are:
\begin{enumerate}
\item A testbed in which tabular Q-learning with a reverse-scramble curriculum covers the full
  state space in about two minutes, so that every network result can be compared with the exact
  optimal policy and value function (\S\ref{sec:testbed}).
\item Four diagnostics that each rule out one explanation of what a value network has learned
  (\S\ref{sec:diagnostics}): a seen/unseen split with a ``steps-to-seen'' statistic, a
  symmetry-image test built from the six whole-cube symmetries that preserve the move set, the
  greedy-versus-beam gap, and a supervised capacity ceiling.
\item An empirical finding on the cube (\S\ref{sec:results}): the network's value function is
  a smooth interpolant in sticker space. It is equally good on seen and unseen states, blind to
  the cube's symmetries, and far more informative than the greedy policy reveals. The
  curriculum stalls because the value error matches the one-move gap between actions, and the
  architecture is not the bottleneck: supervised on the full state space it reaches $79\%$
  greedy success in two minutes, although it can memorise only about a million states exactly.
\end{enumerate}

\section{Related work}

\paragraph{Learning to solve the cube.}
Autodidactic Iteration \citep{mcaleer2018} and DeepCubeA \citep{agostinelli2019} train a value
network on states generated by scrambling the solved cube and solve random instances with a
search (Monte-Carlo tree search, then weighted A$^*$) guided by the network. Both report
success with search; neither separates what the network knows from what the search adds. Our
greedy-versus-beam diagnostic makes that separation on a problem where the optimal answer is
known for every state.

\paragraph{Memorisation and capacity.}
Classical results bound how many random labels a network can fit \citep{cover1965,baum1988};
\citet{zhang2017} showed that modern networks fit random labels on real data. Recent work
measures storage capacity directly: about two bits of knowledge per parameter under
sufficient training \citep{allenzhu2024}, and about $3.6$ bits per parameter for random data in
language models \citep{morris2025}. Our supervised capacity ceiling applies the same idea to a
value network, with the difference that the ``knowledge'' has a known structure, so that the
transition from memorising to generalising can be observed against exact ground truth.

\paragraph{Generalisation beyond capacity.}
Grokking \citep{power2022} and double descent \citep{nakkiran2020} describe regimes in which a
network that could memorise its training set instead, or eventually, finds a generalising
solution. We observe the complementary regime: when the training set exceeds what the network
can memorise, the forced compression yields a policy that solves most random states although
it fits only $56\%$ of the training labels.

\paragraph{Optimisation pathologies in deep RL.}
Value networks trained with bootstrapped targets lose capacity and plasticity over
training \citep{lyle2022,nikishin2022,sokar2023}. Our comparison of Q-learning with supervised
training of the same architecture on the same states isolates the cost of the RL learning
signal from the capacity of the architecture.

\paragraph{Curricula.}
Reverse curricula that start episodes near the goal and move the start distribution
outward \citep{florensa2017} are standard for goal-reaching tasks; we use one and show how a
promotion criterion based on greedy success can lock the curriculum when the value error is of
the order of one move.

\section{Testbed}
\label{sec:testbed}

\paragraph{Environment.}
A state is the colour of each of the 24 stickers. Fixing the down-back-left cubie removes
whole-cube rotations, leaving seven cubies with $7!$ permutations and $3^6$ orientations. The
nine actions are quarter, half and three-quarter turns of the up, right and front layers.
We give every state a perfect-hash index (Lehmer rank of the corner permutation times the
base-3 orientation code), precompute the $3{,}674{,}160\times9$ transition table, and compute
the exact distance $\dist(s)$ of every state by BFS from the solved state. Table~\ref{tab:depth}
gives the distance histogram; $92\%$ of the states, and hence of uniformly random scrambles,
lie at distance 8--10.

\begin{table}[t]
\centering\small
\caption{Number of states at each exact distance from the solved state.}
\label{tab:depth}
\begin{tabular}{lrrrrrrrrrrrr}
\toprule
distance & 0 & 1 & 2 & 3 & 4 & 5 & 6 & 7 & 8 & 9 & 10 & 11 \\
\midrule
states & 1 & 9 & 54 & 321 & 1{,}847 & 9{,}992 & 50{,}136 & 227{,}536 & 870{,}072 & 1{,}887{,}748 & 623{,}800 & 2{,}644 \\
share  & -- & -- & -- & -- & 0.05\% & 0.27\% & 1.4\% & 6.2\% & 23.7\% & 51.4\% & 17.0\% & 0.07\% \\
\bottomrule
\end{tabular}
\end{table}

\paragraph{Learning problem.}
Every move costs $-1$ and an episode ends at the solved state, so with $\gamma=1$ the optimal
action value is $\Q^*(s,a) = -\dist(s')$ for the successor $s'$, and $\max_a \Q^*(s,a) = -\dist(s)$.
The value error of a learned $\Q$ on a state is therefore $|\max_a \Q(s,a) + \dist(s)|$, in moves.
Episodes start from the solved state scrambled by $k\sim\mathrm{Uniform}\{1,\dots,K\}$ random
moves (consecutive moves on different layers). The curriculum depth $K$ starts at~1 and is
promoted when the greedy policy solves at least a threshold fraction of freshly sampled states
at exact distance~$K$; the network agent also trains on starts up to $K+2$.

\paragraph{Tabular agent.}
The table has one row per state. With learning rate $\alpha=1$ (the environment is
deterministic), pessimistic initialisation at $-12$, an ``unknown-action-first'' behaviour
policy and 8{,}192 parallel environments, the table reaches $100\%$ greedy success on random
states, covers every state and drives the mean value error below $0.01$ in about two minutes
($1.8\times10^8$ environment steps). This is the reference for all network experiments.

\paragraph{Network agent.}
The network is a multilayer perceptron that maps the 144-dimensional one-hot encoding of the 24
stickers to nine action values. The main run uses hidden layers 1024--1024--512
($1{,}727{,}497$ parameters), Adam with a cosine learning-rate schedule from $10^{-3}$ to
$10^{-4}$, batches of 4{,}096 states, a target network refreshed every 200 updates, and
all-actions targets $y(s,a) = -1 + \max_{a'} \Q^-(T(s,a),a')$ for every action of every sampled
state \citep{agostinelli2019}. Promotion threshold $0.9$, one hour on an Apple M1~Max.
A second run on an RTX~4090 uses hidden layers 2048--2048--1024 ($6.6$M parameters) for
$5.3$~hours with threshold $0.97$. The main run made $35{,}594$ gradient updates in 63 minutes
($9.4$ updates/s, $2.3\times10^6$ environment steps, $6.5\times10^5$ episodes): $1.5\times10^8$
state rows in total, about 130 per sampled state on average, heavily skewed towards the shallow
states that every episode passes through. \todo{update count and throughput of the 4090 run.}

\section{Diagnostics}
\label{sec:diagnostics}

All diagnostics use the exact distance table and, where needed, the bitmap of states sampled
during training (a state counts as sampled when it entered at least one update). Success means
the greedy policy reaches the solved state within 30 moves.

\paragraph{D1: seen/unseen split and steps-to-seen.}
For each exact distance we sample states and report success separately for sampled and
never-sampled states. If the network memorises, sampled states should be solved markedly more
often. For every solved unseen state we also record how many greedy moves it takes before the
path first enters a sampled state; a median of one or two means success on unseen states is
mostly a step into memorised territory.

\paragraph{D2: symmetry images.}
With one cubie fixed, six whole-cube symmetries preserve the move set: the rotations about the
diagonal through the fixed cubie ($U\!\to\!R\!\to\!F\!\to\!U$) and their compositions with a
mirror that swaps $R$ and $F$ and reverses every turn. Each symmetry $\sigma$ induces a
permutation $\pi$ of the nine moves and a map on states with
$\sigma(T(s,m)) = T(\sigma(s),\pi(m))$ and $\sigma(\text{solved})=\text{solved}$. We construct
$\sigma$ by BFS over the transition table and verify that it is an automorphism, so
$\dist(\sigma(s)) = \dist(s)$ for every state. The image of a sampled state is a state the
network never saw whose solution is the same solution with the moves renamed. If the network
has learned the symmetry, images are solved as often as their originals; if it has memorised,
images are solved as rarely as random unseen states at the same distance.

\paragraph{D3: greedy versus beam.}
The greedy policy commits to $\arg\max_a \Q(s,a)$ at every step. Beam search of width $w$ keeps
the $w$ successor states with the highest $\max_a \Q$ at every depth, with the same network and
no other information. The gap between the two measures how much of the value function's
information the greedy policy discards: a large gap means the ranking of states is right in
aggregate but not reliable one move at a time.

\paragraph{D4: supervised capacity ceiling.}
We remove reinforcement learning from the picture and train the same architecture with the
same optimiser to predict an optimal action for $n$ uniformly random states, with the loss
$-\log\sum_{a\in A^*(s)} p(a\mid s)$ over the set $A^*(s)$ of optimal actions (on average
$2.28$ per state, so the chance level is $25.4\%$). Each $n$ gets the same budget of 6{,}000 Adam
steps and stops early at $99.9\%$ fit. The largest $n$ fitted exactly is the architecture's
memorisation ceiling; held-out accuracy and the greedy solve rate on random states show what the
fit is worth as a policy.

\section{Results}
\label{sec:results}

\subsection{The network run}

The main run's curriculum stalled at $K=7$: greedy success at distance~7 stayed near $52\%$
against a threshold of $90\%$. Training sampled $29.9\%$ of all states. The larger 4090 run also
stalled at $K=7$ with $78.8\%$ of states sampled, $30\%$ greedy success on random states and
$12.6\%$ on never-sampled states. Table~\ref{tab:seen} gives the seen/unseen split of the main
run; measured directly with 2{,}000 states per distance, greedy success on uniformly random
states is $15.8\%$ and the mean value error is $1.07$.

\begin{table}[t]
\centering\small
\caption{D1 on the main run: greedy success and value error for sampled and never-sampled
states at each exact distance (up to 4{,}000 states per distance). Distances 1--7 were sampled
completely. Steps-to-seen: median number of greedy moves from a solved unseen state until the
path enters a sampled state, and the share of solved unseen states needing three or more.}
\label{tab:seen}
\begin{tabular}{rrrrrrrr}
\toprule
& \multicolumn{2}{c}{sampled} & \multicolumn{2}{c}{never sampled} & \multicolumn{2}{c}{steps-to-seen} & value error \\
\cmidrule(lr){2-3}\cmidrule(lr){4-5}\cmidrule(lr){6-7}
distance & share & success & share & success & median & $\ge3$ & sampled / unseen \\
\midrule
1--5 & 100\% & 100\% & -- & -- & -- & -- & 0.01--0.09 \\
6  & 100\%  & 98.9\% & -- & -- & -- & -- & 0.43 \\
7  & 100\%  & 52.3\% & -- & -- & -- & -- & 0.95 \\
8  & 74\%   & 22.4\% & 26\%  & 17.3\% & 1 & 14\% & 0.60 / 0.58 \\
9  & 8.7\%  & 8.4\%  & 91\%  & 9.9\%  & 2 & 22\% & 1.12 / 1.00 \\
10 & 0.75\% & 0.0\%$^\dagger$ & 99\% & 5.9\% & 3 & 58\% & 3.66 / 2.05 \\
11 & 5.7\%  & 0.0\%  & 94\%  & 1.6\%  & 4 & 97\% & 5.04 / 3.35 \\
\bottomrule
\multicolumn{8}{l}{\footnotesize $^\dagger$30 states.}
\end{tabular}
\end{table}

\subsection{D1: seen and unseen states are solved equally often}

At every distance where both groups exist, sampled and never-sampled states are solved at
nearly the same rate ($22.4\%$ vs.\ $17.3\%$ at distance~8, $8.4\%$ vs.\ $9.9\%$ at distance~9)
with nearly the same value error. Being sampled confers no advantage at depth. At distance~10
the sampled states are solved \emph{less} often than unseen ones. Distance-10 states are never
episode starts; the ones that get sampled are those an episode walked into from distance~9,
that is, states the network misjudged as closer to the goal. The sampled set at depth is
therefore a biased sample of the network's own errors.

The steps-to-seen statistic shows that most solved unseen states at distances 8--9 enter
sampled territory within one or two moves. This does not indicate memorisation, since $74\%$ of
distance-8 states are sampled and a step into that set is unavoidable; at distance~10, $58\%$ of
the solved unseen states take three or more moves before touching any sampled state.

\subsection{D2: the network does not know the cube's symmetries}

\begin{table}[t]
\centering\small
\caption{D2 on the main run: greedy success on sampled states, on their never-sampled symmetric
images, and on random never-sampled states, 2{,}000 states per cell. Distances 1--7 have no
unseen images because every state was sampled.}
\label{tab:sym}
\begin{tabular}{rrrrr}
\toprule
distance & sampled & symmetric images (unseen) & random unseen & image / random \\
\midrule
8  & 22.9\% & 18.1\% & 19.1\% & 1.0 \\
9  & 12.0\% & 9.6\%  & 9.7\%  & 1.0 \\
10 & 0.4\%  & 1.7\%  & 6.0\%  & 0.3 \\
11 & 0.0\%$^\ddagger$ & 0.0\%$^\ddagger$ & 2.1\% & -- \\
\bottomrule
\multicolumn{5}{l}{\footnotesize $^\ddagger$158 and 103 states.}
\end{tabular}
\end{table}

All five non-identity symmetries were verified as automorphisms of the transition graph.
Table~\ref{tab:sym} shows that the symmetric images of sampled states are solved exactly as
often as random unseen states at the same distance. A solution the network can produce for a
state does not transfer to the same solution with the layers renamed. Since the one-hot sticker
encoding does not make the symmetry explicit, a multilayer perceptron has no reason to become
equivariant, and it did not.

Together, D1 and D2 leave neither of the two usual explanations standing: the network does not
store individual states, and it has not learned the group structure.

\subsection{D3: the value function knows far more than the greedy policy shows}

\begin{table}[t]
\centering\small
\caption{D3 on the main run: greedy versus beam search (width 32) with the same network.
Greedy figures at distances 8--11 are the sampled/unseen averages of Table~\ref{tab:seen}.}
\label{tab:beam}
\begin{tabular}{rrr}
\toprule
distance & greedy & beam, width 32 \\
\midrule
$\le6$ & $\ge 98.9\%$ & 100\% \\
7  & 52\%  & 100\% \\
8  & 21\%  & 100\% \\
9  & 10\%  & 92.5\% \\
10 & 6\%   & 66.0\% \\
11 & 2\%   & 40.5\% \\
\midrule
uniformly random state & $\approx16\%$ & 90.3\% \\
\bottomrule
\end{tabular}
\end{table}

Beam search of width~32 solves $90.3\%$ of uniformly random states, against about $16\%$ for
the greedy policy, and $100\%$ of distance-8 states although a quarter of them were never
sampled and distances 10--11 were never episode starts (Table~\ref{tab:beam}). The value
function ranks states correctly in aggregate throughout the state space, including the $70\%$ it
never sampled. What it lacks is precision: the true values of adjacent actions differ by
exactly one move, the value error at distance~7 is $0.95$ moves, and a greedy choice between
two actions whose estimates differ by less than the error is a coin flip. Beam search tolerates
individual misrankings; greedy compounds them.

\subsection{Why the curriculum stalls}

The promotion criterion asks the greedy policy to solve $90\%$ of distance-7 states. Every
distance-7 state was sampled and beam search solves all of them, so the network ``knows'' the
region; greedy success there is $52\%$ because the value error equals the signal. The
curriculum therefore never advances, distances 10--11 never become episode starts, and their
values are extrapolated: the value error grows to $2$--$3.4$ moves and the network
systematically underestimates the distance of deep states (Table~\ref{tab:seen}). The stall is
self-locking, and it is a property of the promotion criterion, not of the network.

\subsection{D4: capacity is not the bottleneck}

\begin{table}[t]
\centering\small
\caption{D4: the main run's architecture (1.73M parameters, $\approx3.5$ Mbit at two bits per
parameter; a full policy is $3{,}674{,}160\times\log_2 9 \approx 11.6$ Mbit) trained supervised on
$n$ random states with 6{,}000 Adam steps each. Chance level for accuracy is $25.4\%$.}
\label{tab:cap}
\begin{tabular}{rrrrr}
\toprule
$n$ (states) & steps & fit & held-out & greedy solve, random states \\
\midrule
300{,}000   & 1{,}500 & 100.0\% & 29.4\% & 0.0\% \\
1{,}000{,}000 & 6{,}000 & 92.2\% & 32.1\% & 4.5\% \\
2{,}000{,}000 & 6{,}000 & 64.5\% & 39.8\% & 42.3\% \\
3{,}674{,}160 & 6{,}000 & 56.1\% & -- & 79.2\% \\
\bottomrule
\end{tabular}
\end{table}

Table~\ref{tab:cap} shows three things. First, the architecture can memorise, but only about
$0.3$--$1$ million states: $300{,}000$ states are fitted exactly in 1{,}500 steps, and the
fitted information ($\approx1$--$3$ Mbit) is consistent with the two-bits-per-parameter
rule \citep{allenzhu2024}. What is memorised does not transfer: held-out accuracy at
$n=300{,}000$ is at chance. Second, once $n$ exceeds the ceiling the network is forced to
compress and starts to generalise: at $n=2$ million the fit drops to $64.5\%$ while held-out
accuracy rises to $39.8\%$ and greedy solve rate jumps to $42\%$; on the full state space a
per-state accuracy of $56\%$ yields $79\%$ greedy success, because most ``wrong'' actions keep
the distance rather than increase it. Third, the same architecture that Q-learning brought to
about $16\%$ greedy success in an hour reaches $79\%$ in two minutes when given clean targets.
The bottleneck of the RL run is the learning signal, bootstrapped from the network's own
noisy and non-stationary estimates, not the parameter count.

\subsection{What a larger network buys}

The 4090 run with four times the parameters raised greedy success on random states from about
$16\%$ to $30\%$ and stalled at the same $K=7$. In the light of D3 and D4, parameters buy
\emph{precision} of the value estimates: a smaller error lets the greedy policy survive a few
more moves, but as long as the error is of the order of the one-move gap the promotion criterion
is not met. Parameters do not buy per-state memory, which D1 shows the network does not use.

\section{Discussion}

\paragraph{What the network learned.}
The three explanations in the introduction are each contradicted by one diagnostic, leaving a
fourth: the value function is a smooth interpolant in the sticker representation. States that
share many stickers receive similar values. This is enough for beam search to follow the
gradient of the true distance across the whole state space, and not enough for one-move
decisions near the resolution limit of the estimate. The network is neither a compressed table
nor a solver; it is a heuristic of roughly one move's accuracy.

\paragraph{Consequences for method design.}
The promotion criterion of a reverse curriculum should not depend on the greedy policy when the
value error is of the order of the action gap; one-step accuracy or search success are
criteria that advance. If structural knowledge is wanted, the input representation, not the
network size, is the lever: an encoding of cubie positions and orientations, or symmetry
augmentation of every sample, makes the symmetries learnable; a wider MLP on stickers does not.

\paragraph{Limitations.}
Every configuration was run with a single seed. The main and 4090 runs differ in size, budget
and threshold, so the comparison between them is qualitative.
All results are on one puzzle; the diagnostics are general, the conclusions are not yet.

\paragraph{Planned experiments.}
Three seeds per configuration; a width sweep from 256 to 4096 at a fixed number of updates with
per-size learning rates; the supervised sweep over parameters and $n$ to map the boundary
between memorising and generalising; a second exhaustively solvable domain; and the curriculum
with a one-step or search-based promotion criterion, resumed from the stalled checkpoint.

\section{Conclusion}

On a puzzle where every state's optimal value is known, four cheap diagnostics show that a
value network trained by Q-learning neither memorises its training states nor learns the
puzzle's symmetries, that its value function is nevertheless informative over the entire state
space, and that its failure to advance is a matter of precision and learning signal rather than
capacity. Problems of this kind, exhaustively solvable yet too large for a network to store, are
a useful place to ask what deep RL representations contain before scaling to problems where the
question cannot be answered.

\section*{Reproducibility}
Code, the training monitor and the diagnostic scripts (\texttt{eval --seen}, \texttt{eval
--beam}, \texttt{symmetry\_test.py}, \texttt{capacity.py}) are available at
\url{https://github.com/u112358/RLDemo}. All runs are reproducible from the command lines in the appendix.

\bibliographystyle{plainnat}
\begin{thebibliography}{99}
\providecommand{\natexlab}[1]{#1}
% every entry verified against arXiv / PMLR / OpenReview on 2026-09-15
\bibitem[Agostinelli et~al.(2019)]{agostinelli2019}
F.~Agostinelli, S.~McAleer, A.~Shmakov, and P.~Baldi.
\newblock Solving the Rubik's cube with deep reinforcement learning and search.
\newblock \emph{Nature Machine Intelligence}, 1:356--363, 2019.

\bibitem[Allen-Zhu and Li(2024)]{allenzhu2024}
Z.~Allen-Zhu and Y.~Li.
\newblock Physics of language models: Part 3.3, knowledge capacity scaling laws.
\newblock \emph{arXiv:2404.05405}, 2024.

\bibitem[Baum(1988)]{baum1988}
E.~B. Baum.
\newblock On the capabilities of multilayer perceptrons.
\newblock \emph{Journal of Complexity}, 4(3):193--215, 1988.

\bibitem[Cover(1965)]{cover1965}
T.~M. Cover.
\newblock Geometrical and statistical properties of systems of linear inequalities with applications in pattern recognition.
\newblock \emph{IEEE Transactions on Electronic Computers}, EC-14(3):326--334, 1965.

\bibitem[Florensa et~al.(2017)]{florensa2017}
C.~Florensa, D.~Held, M.~Wulfmeier, M.~Zhang, and P.~Abbeel.
\newblock Reverse curriculum generation for reinforcement learning.
\newblock In \emph{Conference on Robot Learning}, PMLR 78:482--495, 2017.

\bibitem[Lyle et~al.(2022)]{lyle2022}
C.~Lyle, M.~Rowland, and W.~Dabney.
\newblock Understanding and preventing capacity loss in reinforcement learning.
\newblock In \emph{International Conference on Learning Representations}, 2022.

\bibitem[McAleer et~al.(2019)]{mcaleer2018}
S.~McAleer, F.~Agostinelli, A.~Shmakov, and P.~Baldi.
\newblock Solving the Rubik's cube with approximate policy iteration.
\newblock In \emph{International Conference on Learning Representations}, 2019. (Preprint: Solving the Rubik's cube without human knowledge, arXiv:1805.07470, 2018.)

\bibitem[Morris et~al.(2025)]{morris2025}
J.~X. Morris, C.~Sitawarin, C.~Guo, N.~Kokhlikyan, G.~E. Suh, A.~M. Rush, K.~Chaudhuri, and S.~Mahloujifar.
\newblock How much do language models memorize?
\newblock \emph{arXiv:2505.24832}, 2025.

\bibitem[Nakkiran et~al.(2020)]{nakkiran2020}
P.~Nakkiran, G.~Kaplun, Y.~Bansal, T.~Yang, B.~Barak, and I.~Sutskever.
\newblock Deep double descent: Where bigger models and more data hurt.
\newblock In \emph{International Conference on Learning Representations}, 2020.

\bibitem[Nikishin et~al.(2022)]{nikishin2022}
E.~Nikishin, M.~Schwarzer, P.~D'Oro, P.-L. Bacon, and A.~Courville.
\newblock The primacy bias in deep reinforcement learning.
\newblock In \emph{International Conference on Machine Learning}, PMLR 162:16828--16847, 2022.

\bibitem[Power et~al.(2022)]{power2022}
A.~Power, Y.~Burda, H.~Edwards, I.~Babuschkin, and V.~Misra.
\newblock Grokking: Generalization beyond overfitting on small algorithmic datasets.
\newblock \emph{arXiv:2201.02177}, 2022.

\bibitem[Sokar et~al.(2023)]{sokar2023}
G.~Sokar, R.~Agarwal, P.~S. Castro, and U.~Evci.
\newblock The dormant neuron phenomenon in deep reinforcement learning.
\newblock In \emph{International Conference on Machine Learning}, PMLR 202:32145--32168, 2023.

\bibitem[Zhang et~al.(2017)]{zhang2017}
C.~Zhang, S.~Bengio, M.~Hardt, B.~Recht, and O.~Vinyals.
\newblock Understanding deep learning requires rethinking generalization.
\newblock In \emph{International Conference on Learning Representations}, 2017.
\end{thebibliography}

\appendix
\section{Hyperparameters and commands}

\paragraph{Tabular agent.}
\texttt{python play\_rubik.py train}: 8{,}192 environments, $\alpha=1$, $\gamma=1$,
$\epsilon=0.1$, initial value $-12$ with clamp, promotion threshold $0.97$, episode cap 25,
evaluation every 50 iterations. Reaches $K=11$ and $\ge99.9\%$ success and coverage in
$1.8\times10^8$ environment steps, 127~s on the reference machine.

\paragraph{Network agent, main run.}
\texttt{python play\_rubik.py train --agent net --backend torch --tag main --batch 4096
--minutes 60 --lr-final 1e-4 --promote 0.9}: layers 1024--1024--512, 64 parallel
environments, all-actions targets, target network every 200 updates, curriculum margin 2,
$\epsilon=0.1$, episode cap 25, evaluation every 50 updates. Apple M1~Max, PyTorch MPS.

\paragraph{Diagnostics.}
\texttt{eval --agent net --tag main --seen}, \texttt{eval --agent net --tag main --beam 32},
\texttt{python symmetry\_test.py --tag main}, \texttt{python capacity.py --layers 1024,1024,512
--n 300000,1000000,2000000,3674160 --steps 6000 --batch 8192}.

\section{Constructing the symmetries}
Let $\pi_\rho$ map move $(\ell, t)$ (layer $\ell\in\{U,R,F\}$, turn $t\in\{1,2,3\}$ quarter
turns) to $(\ell+1 \bmod 3, t)$, and $\pi_\mu$ map it to $(\mathrm{swap}_{RF}(\ell), 4-t)$.
Starting from $\sigma(\text{solved})=\text{solved}$, BFS over the transition table assigns
$\sigma(T(s,m)) := T(\sigma(s),\pi(m))$ and checks every already-assigned value for consistency.
All five non-identity elements of the group generated by $\pi_\rho$ and $\pi_\mu$ passed the
check and preserve the BFS distance of every state.

\end{document}
